For fixed $k\ge2$, the question asks for an asymptotic mean of the least $k$th-power nonresidue. Elliott's original convention is essential: define $n_k(p)$ as that least positive nonresidue when $p\equiv1\pmod k$, and put $n_k(p)=0$ for all other primes. This is explicitly the first paragraph of \emph{A problem of Erd\H{o}s concerning power residue sums}, Acta Arithmetica 13 (1967), \url{https://doi.org/10.4064/aa-13-2-131-149}. Without a convention, the unrestricted definition can fail even to exist, since if $\gcd(k,p-1)=1$, every nonzero residue is a $k$th power. We use Elliott's convention throughout.

The known result is affirmative. This attempt proves all fixed-value densities, constructs the finite positive candidate constant, and proves the corresponding lower limit. It explicitly leaves the uniform prime-average tail estimate unproved, so does not claim a complete reconstruction of Elliott's analytic theorem.
